\documentclass{cv}

\usepackage{fontspec}
\usepackage[british]{babel}
\usepackage[left=1.0cm,top=1.0cm,right=1.0cm,bottom=2.0cm]{geometry}

\usepackage[activate={true, nocompatibility}, final, tracking=true, kerning=false, spacing=false, factor=1050, stretch=10, shrink=10]{microtype}
\microtypecontext{spacing=nonfrench}

\tolerance=1414
\hbadness=1414
\emergencystretch=1.5em
\sloppy

\usepackage{setspace}

\usepackage{tabularx}

\usepackage[inline]{enumitem}

\usepackage{needspace}
\usepackage{fancyhdr}
\usepackage{hyperref}

\hypersetup{
    hidelinks,
    pdftitle={Curriculum Vitae - Alexander C. Whitehead},
    pdfauthor={Alexander C. Whitehead}
}

\setlength{\footskip}{1.0cm}
\pagestyle{fancy}
\fancyhf{}
\fancyfoot[C]{\footnotesize Dr. Alexander C. Whitehead \textbar{} Curriculum Vitae \textbar{} \thepage}
\renewcommand{\headrulewidth}{0pt}
\renewcommand{\footrulewidth}{0pt}

\setmainfont[
    Path=./,
    UprightFont = {*-Regular},
    BoldFont    = {*-Bold},
    ItalicFont  = {*-Italic},
    BoldItalicFont = {*-BoldItalic}
]{DINish.otf}

\makeatletter

\def\@photo{}
\def\@photowidth{0pt}
\newcommand{\photo}[2]{\def\@photo{#1}\def\@photowidth{#2}}

\makeatother

\begin{document}
    \setlength{\parindent}{0pt}

    {\LARGE\textbf{Dr.~Alexander~C.~Whitehead}}\\
    {\small
    Fuhlsbüttler Straße 599, 22337 Hamburg, Germany\\
    \href{tel:+4915563707429}{+49 15563 707429} \textbar{} \href{mailto:alexandercwhitehead@gmail.com}{alexandercwhitehead@gmail.com}\\
    \href{https://scholar.google.com/citations?user=-hH_J28AAAAJ}{Google Scholar}
    \textbar{}
    \href{https://orcid.org/0000-0002-9151-3295}{ORCID}
    \textbar{}
    \href{https://github.com/ALEXJAZZ008008}{GitHub}\\
    \textbf{Authorised to work in Germany; no visa sponsorship required \textbar{} Available immediately}}
    
    \begin{spacing}{1.1}
        \normalsize
        
        \vspace{0.35cm}
        \begin{rSection}{Profile}
            Scientific software engineer and applied machine learning scientist with a PhD in medical imaging and BSc/MSc degrees in computer science. I design, implement and evaluate scientific software, machine learning models and quantitative imaging algorithms, primarily in Python, drawing on substantial earlier C++ experience. I independently developed ASTRA for non-invasive PET quantification and collaborated with GE HealthCare on PET/CT respiratory motion correction.
        \end{rSection}
        
        \begin{rSection}{Technical Skills}
            \begin{itemize}
                \item \textbf{Python (PhD--present):} NumPy, SciPy, pandas, Matplotlib, Jupyter, scikit-learn, TensorFlow/Keras, PyTorch; machine learning, data analysis and scientific computing.
                \item \textbf{C++ (BSc--mid-PhD):} DirectX, OpenGL, Qt, PCL, CMake; graphics, simulation and scientific software.
                \item \textbf{Machine Learning \& Statistics:} Deep learning (CNNs, autoencoders, transformers, diffusion models, Bayesian neural networks), fine-tuning, learned embeddings, k-NN and cosine similarity, interpretability; PCA, weighted linear regression, survival analysis, uncertainty quantification, cross-validation.
                \item \textbf{Other Programming Languages:} C (BSc; compilers, embedded systems), C\#/.NET (primary BSc language; Unity, Xamarin), Java (BSc; Android), JavaScript (BSc), MATLAB (PhD), SQL (MSc).
                \item \textbf{Systems \& Tools:} Linux (primary operating system), Git/GitHub (pull requests, collaborative development), Bash, HPC clusters (SGE, SLURM).
                \item \textbf{Medical Imaging \& 3D Sensing:} PET/CT, dynamic PET, image reconstruction and registration, respiratory motion correction, kinetic modelling; LiDAR and Kinect depth sensing, point-cloud registration and processing.
            \end{itemize}
        \end{rSection}
        
        \begin{rSection}{Experience}
            \vspace{-0.35cm}
            
            \item \begin{rSubsection}{London South Bank University \textnormal{-- London, UK}}{Feb 2025 -- Present}{Visiting Fellow}{}
                \item \begin{itemize}
                    \item Independently designed and implemented ASTRA in TensorFlow, integrating a spatiotemporal patched autoencoder, transformer and metabolite-correction module to estimate metabolite-corrected arterial input functions with predictive uncertainty from dynamic [\textsuperscript{11}C]PBR28 PET without arterial blood sampling at inference.
                    \item Co-first-authored the \href{https://doi.org/10.21203/rs.3.rs-8895728/v1}{2026 ASTRA preprint}, reporting evaluation across four PET datasets and better agreement with blood-derived quantification than a population-corrected image-derived input-function baseline on two datasets.
                    \item Co-first-authored a \href{https://doi.org/10.1109/EMBC58623.2025.11251612}{2025 IEEE EMBC paper} on multicentre neural-network PET quantification.
                \end{itemize}
            \end{rSubsection}
        
            \item \begin{rSubsection}{University College London \textnormal{-- London, UK}}{Sep 2018 -- Sep 2025}{Doctoral Researcher}{}        
                \item \begin{itemize}
                    \item Designed, implemented and evaluated PET/CT respiratory motion-correction algorithms using image registration and \href{https://doi.org/10.1109/NSS/MIC44845.2022.10399154}{motion models fitted to coarsely gated data and applied to finer gates}, improving image contrast and quantitative accuracy in simulations funded by GE HealthCare.
                    \item Developed a data-driven method for retrospective respiratory signal extraction from dynamic PET data without external monitoring equipment.
                    \item Developed reusable Python and MATLAB pipelines for PET simulation, image reconstruction and quantitative evaluation of respiratory motion-correction methods.
                    \item First-authored a \href{https://doi.org/10.1109/NSSMICRTSD49126.2023.10338687}{2023 IEEE NSS/MIC paper} on Bayesian neural networks for PET input-function estimation with uncertainty quantification.
                    \item First-authored a \href{https://doi.org/10.1088/1361-6560/ad5ef1}{2024 \textit{Physics in Medicine \& Biology} paper} (Editor's Choice), demonstrating earlier and more accurate respiratory signal extraction using selective PCA than conventional PCA on seven held-out acquisitions from a 22-acquisition study.
                    \item Contributed code to \href{https://github.com/UCL/STIR}{STIR}, including \href{https://github.com/UCL/STIR/blob/a56de48ecb020547ea40ccb992e956757a8c5006/src/include/stir/numerics/FastErf.inl}{C++ error-function interpolation}, and to \href{https://github.com/SyneRBI/SIRF}{SIRF}, both open-source image-reconstruction frameworks.
                    \item Collaborated with GE HealthCare engineers and hospital staff on PET/CT respiratory motion-correction research.
                    \item Supervised four MSc research projects, serving as primary supervisor for one.
                \end{itemize}
            \end{rSubsection}
        
            \item \begin{rSubsection}{University College London \textnormal{-- London, UK}}{Dec 2022 -- Dec 2024}{Research Fellow}{}
                \item \begin{itemize}
                    \item Developed CT-based models for survival and disease progression in idiopathic pulmonary fibrosis, including Grad-CAM interpretability analysis.
                    \item Evaluated survival models on 550 CT acquisitions using five-fold cross-validation; likelihood-based models outperformed the Cox-based approaches tested.
                    \item Built reproducible Python pipelines for survival-model evaluation.
                    \item First-authored an \href{https://doi.org/10.1109/NSSMICRTSD49126.2023.10337922}{IEEE NSS/MIC paper} on CT-based survival analysis.
                    \item Co-authored \href{https://doi.org/10.1016/j.media.2023.103016}{\textit{CenTime}} on event-conditional censoring, published in \textit{Medical Image Analysis}.
                \end{itemize}
            \end{rSubsection}
        
            \item \begin{rSubsection}{University College London \textnormal{-- London, UK}}{Oct 2018 -- Dec 2024}{Senior Postgraduate Teaching Assistant}{}
                \item \begin{itemize}
                    \item Lectured on introductory programming to approximately 200 students.
                    \item Taught electronics, embedded systems programming in C, engineering mathematics using MATLAB, Python programming, medical physics and professional writing.
                \end{itemize}
            \end{rSubsection}
        
            \item \begin{rSubsection}{University of Hull \textnormal{-- Hull, UK}}{Oct 2017 -- Sep 2018}{Demonstrator}{}
                \item \begin{itemize}
                    \item Supported undergraduate teaching in C and C\# programming, electronics, embedded systems and databases.
                \end{itemize}
            \end{rSubsection}
        
            \item \begin{rSubsection}{University of Hull (Digital Centre) \textnormal{-- Hull, UK}}{Jun 2017 -- Sep 2017}{Research Intern}{}
                \item \begin{itemize}
                    \item Designed and built a drone-mounted rotating-LiDAR scanner for 3D room mapping; contributed to an experimental autonomous-vehicle project.
                \end{itemize}
            \end{rSubsection}
        \end{rSection}
        
        \begin{rSection}{Education}
            \vspace{-0.35cm}
            
            \item \begin{rSubsection}{University College London \textnormal{-- London, UK}}{2025}{PhD in Medical Imaging Physics and Medical Image Computing}{}
                \item \textbf{Thesis:} \href{https://discovery.ucl.ac.uk/id/eprint/10212970/}{\textit{Improved Quantification for Respiratory Gated PET/CT: Data-Driven Algorithms for Respiratory Motion Correction in PET/CT}}
            \end{rSubsection}
        
            \item \begin{rSubsection}{University of Hull \textnormal{-- Hull, UK}}{2018}{MSc in Advanced Computer Science}{Distinction}
                \item \begin{itemize}
                    \item Developed a C++/Qt framework for Kinect-based point-cloud registration and motion tracking.
                    \item Implemented a DirectX 11 graphics engine with custom HLSL shaders, lighting, shadows and physics.
                \end{itemize}
            \end{rSubsection}
        
            \item \begin{rSubsection}{University of Hull \textnormal{-- Hull, UK}}{2017}{BSc in Computer Science}{First Class Honours}
                \item \begin{itemize}
                    \item Developed a C\#/Xamarin GPS-based multiplayer game using TCP/UDP networking and geometric algorithms.
                    \item Implemented a compiler in C for a custom language using Flex/Bison and constant-folding optimisation.
                \end{itemize}
            \end{rSubsection}
        \end{rSection}
        
        \begin{rSection}{Languages}
            \textbf{English:} Native \hspace{1cm}\textbar{}\hspace{1cm} \textbf{German:} A1 (certified; continuing study)
        \end{rSection}
        
        \begin{rSection}{Personal Interests}
            Gardening is my main hobby; I tend indoor and outdoor plants and plan my tomato-growing season using grow lights. I also design and track strength-training programmes, and enjoy live music and fish-keeping.
        \end{rSection}
    \end{spacing}
\end{document}
